\section{Purpose of laboratory work}

The purpose of this laboratory work is to implement an \textbf{actuator control system}
on the ESP32 microcontroller using FreeRTOS for concurrent task execution. The system
controls two actuators -- an analog actuator (DC motor via L298N PWM) and a binary
actuator (relay module) -- with commands received over serial, a signal conditioning
pipeline, and real-time reporting on an OLED display.
This is Variant~C (100\%), requiring both analog and binary actuator control.

\textbf{Hardware note:} During initial development and testing, an LED connected to
GPIO~18 (ENA pin of the L298N) with a $220\,\Omega$ resistor was used in place of a
DC motor, as the motor was not yet available. The LED brightness varies proportionally
with the PWM duty cycle, allowing the ramp-up, ramp-down, and speed control behavior
to be visually verified. During the actual laboratory presentation, the system was
demonstrated with a real \textbf{DC motor connected through the L298N H-bridge driver},
which responded correctly to all speed, direction, and emergency stop commands.
The L298N direction pins (GPIO~16, GPIO~17) were wired and driven correctly throughout.
The driver abstraction design ensured that no code changes were required when
switching from the LED to the motor.

\subsection{Objectives of the laboratory work}

Laboratory work has the following objectives that should be achieved:

\begin{itemize}
  \item Implementing a serial command interface for receiving actuator commands
    (\texttt{S<0-100>}, \texttt{DF}, \texttt{DR}, \texttt{X}, \texttt{RON},
    \texttt{ROFF}, \texttt{?}) at a 50\,ms polling period
  \item Implementing analog actuator control through a four-stage command conditioning
    pipeline: saturation (clamping to $0$--$100\,\%$), a median filter (window size~5)
    for impulsive command noise removal, an exponential moving average
    (EMA, $\alpha = 0.3$) for smoothing, and a ramping stage ($2\,\%$ per 50\,ms cycle)
    that limits the rate of change applied to the motor
  \item Controlling a DC motor (L298N, GPIO~16/17/18) via PWM (LEDC 5\,kHz, 8-bit)
    with direction control (forward/reverse) and an emergency stop that bypasses the
    ramp and resets all filters
  \item Controlling a single-channel relay module (GPIO~19, active-HIGH) as a binary
    actuator with ON/OFF target tracking
  \item Evaluating alert conditions (OVERLOAD $\geq 95\,\%$, LIMIT\_REACHED at
    $0\,\%$ or $100\,\%$) and driving LED indicators accordingly
  \item Displaying a structured report at 500\,ms on an OLED 128$\times$64 display
    showing motor state, direction, target/conditioned/current speed, relay state,
    and alert; emitting a serial plotter line every 500\,ms
  \item Structuring the firmware into four independent FreeRTOS tasks (Command,
    Conditioning, Actuator, Reporter) communicating through a binary semaphore and a mutex
  \item Organizing the codebase into three hardware-abstraction layers: MCAL, ECAL,
    and SRV
\end{itemize}

\subsection{Problem definition}

The application is structured in four distinct FreeRTOS tasks:

\begin{enumerate}
  \item \textbf{Task 1 -- Command (TaskCommand):} Non-blocking serial read every 50\,ms.
    Parses commands, validates ranges, stores target speed, direction, relay target, and
    emergency stop flag in shared \texttt{ActuatorData}. Signals the conditioning task via
    a binary semaphore after each valid speed command.

  \item \textbf{Task 2 -- Signal Conditioning (TaskConditioning):} Runs every 50\,ms,
    blocked on the binary semaphore with timeout. Applies a four-stage pipeline:
    \begin{enumerate}
      \item \textbf{Saturation} -- clamps the target to $0$--$100\,\%$
      \item \textbf{Median filter} (window = 5) -- removes impulsive command noise
      \item \textbf{Exponential Moving Average} ($\alpha = 0.3$) -- smooths fluctuations:
        $y_n = \alpha \cdot x_n + (1 - \alpha) \cdot y_{n-1}$
      \item \textbf{Ramping} -- increments \texttt{currentSpeed} toward
        \texttt{conditionedSpeed} by at most $2\,\%$ per cycle ($\approx 2.5\,\mathrm{s}$
        full ramp), manages motor state transitions (IDLE / RAMP\_UP / RUNNING /
        RAMP\_DOWN)
    \end{enumerate}
    Evaluates alerts and drives green/red LED indicators.

  \item \textbf{Task 3 -- Actuator (TaskActuator):} Runs every 50\,ms. Reads
    \texttt{currentSpeed} and \texttt{direction} from \texttt{ActuatorData}, maps
    $0$--$100\,\%$ to PWM duty $0$--$255$ and applies it via the L298N LEDC driver.
    Reads \texttt{relayTarget} and applies it to the relay driver when it differs from
    the current relay state.

  \item \textbf{Task 4 -- Reporter (TaskReporter):} Runs every 500\,ms using
    \texttt{vTaskDelayUntil}. Updates the OLED 128$\times$64 with all actuator state
    fields, emits a serial plotter line in \texttt{>name:value} format, and every
    5\,seconds prints a full structured report to STDIO.
\end{enumerate}
